% =====================================================================
\chapter{Обновление прошивки и SD-карты}
\label{ch:firmware}
% =====================================================================
\chapterintro
  {как узнать версию EdgeTX, сделать резервную копию, обновить прошивку пульта, содержимое
   SD-карты, внутренний модуль ExpressLRS и мультипротокольный модуль}
  {компьютер с браузером Chrome/Edge, кабель USB-C с передачей данных, картридер (желательно)}

\section{Нужно ли обновляться}

Пульты приходят с завода с рабочей версией EdgeTX, и летать можно сразу. Обновляться стоит,
если:
\begin{itemize}
  \item нужна новая функция (например, из заметок к выпуску);
  \item исправлена ошибка, которая вас касается;
  \item версия внутреннего модуля ELRS устарела и не совпадает с приёмниками;
  \item вы собираете единый парк пультов (учебная группа!) и хотите одинаковую версию везде.
\end{itemize}

Правило хорошего тона: \textbf{не обновляйтесь накануне соревнований или важного полёта} и
не ставьте ночные сборки (nightly) на основной пульт.

\section{Узнать текущую версию}

\mpath{SYS\sep About} (последняя страница системного меню, листайте \btn{PAGE>}) показывает
версию прошивки, например \code{edgetx-tx12mk2-v2.11.2}. Там же видно, для какого пульта
собрана прошивка --- это важно: у TX12 MKI, TX12 MKII и Boxer \emph{разные} файлы прошивки.

\begin{caution}
Прошивка от другого пульта может не запуститься или неправильно работать с кнопками и
стиками. При выборе файла внимательно проверяйте название: \code{tx12}, \code{tx12mk2},
\code{boxer}.
\end{caution}

\section{Резервная копия}

\begin{enumerate}
  \item Включите пульт и подключите кабель USB. На экране появится вопрос о режиме USB ---
        выберите \menu{USB Storage (SD)}. Пульт появится на компьютере как флешка.
  \item Скопируйте \emph{всё} содержимое SD-карты в отдельную папку на компьютере
        (например, \code{backup-2026-09-26}). Особенно важны папки \code{MODELS} и \code{RADIO}.
  \item Безопасно извлеките диск и отключите кабель.
\end{enumerate}

Можно также вынуть microSD-карту и скопировать её через картридер.

\section{Обновление через EdgeTX Buddy}

\textbf{EdgeTX Buddy} (\code{buddy.edgetx.org}) --- веб-приложение, которое работает прямо в
браузере на основе Chromium (Chrome, Edge, Яндекс.Браузер) через WebUSB. Существует и
настольная версия.

\subsection{Шаг 1. Прошивка пульта}

\begin{enumerate}
  \item Откройте Buddy, вкладка \menu{Flash} (прошивка).
  \item Выберите версию (последнюю стабильную --- \emph{release}, не \emph{nightly})
        и модель пульта: \menu{RadioMaster TX12 Mark II}, \menu{RadioMaster TX12} или
        \menu{RadioMaster Boxer}.
  \item При желании в дополнительных параметрах выберите язык интерфейса (в т.\,ч. русский)
        и опции сборки.
  \item \textbf{Выключите} пульт и подключите его к компьютеру кабелем. Пульт должен определиться
        как устройство \code{STM32 BOOTLOADER} (режим DFU, экран при этом может оставаться
        тёмным).
  \item Нажмите \menu{Flash} и выберите устройство в окне браузера. Дождитесь окончания
        записи (1--2 минуты). Не отключайте кабель!
\end{enumerate}

\begin{tip}
На Windows устройство DFU иногда не видно в браузере, потому что к нему не привязан драйвер
WinUSB. Buddy сам предложит инструкцию; обычно достаточно один раз выполнить привязку
драйвера утилитой Zadig для устройства \code{STM32 BOOTLOADER}. На Linux может понадобиться
udev-правило для доступа к USB без прав суперпользователя.
\end{tip}

\subsection{Шаг 2. Содержимое SD-карты}

Каждая версия EdgeTX ожидает на SD-карте файлы определённой версии: звуки, скрипты,
служебные картинки. На вкладке \menu{SD Card} в Buddy:
\begin{enumerate}
  \item включите пульт в режиме \menu{USB Storage (SD)};
  \item выберите папку диска пульта, версию пакета и язык голосового пакета;
  \item Buddy запишет нужные файлы, \emph{не трогая} ваши модели и настройки.
\end{enumerate}

Вместо Buddy можно скачать архив \code{sdcard-bw128x64-...zip} (для монохромных экранов
128×64) и голосовой пакет с GitHub и распаковать их на карту поверх старых файлов.

\section{Обновление через загрузчик (bootloader)}

Если DFU недоступен, прошивку можно поставить со SD-карты через встроенный загрузчик EdgeTX.

\begin{enumerate}
  \item Скачайте файл прошивки \code{.bin} для вашего пульта и положите его в папку
        \code{FIRMWARE} на SD-карте.
  \item Выключите пульт. \textbf{Прижмите оба горизонтальных триммера внутрь}
        (к центру пульта) и, удерживая их, включите питание. Появится экран загрузчика.
  \item Выберите \menu{Write Firmware}, затем файл, и подтвердите (обычно долгим \btn{ENTER}).
  \item По окончании выберите \menu{Exit} или перезагрузите пульт.
\end{enumerate}

\begin{figure}[H]
\centering
\begin{tikzpicture}[scale=0.8]
  \fill[black!10,rounded corners=10pt] (-4.4,-1.2) rectangle (4.4,1.2);
  \foreach \x in {-2.8,2.8}{\draw[fill=gray!30,rounded corners=1pt] (\x-0.6,-0.1) rectangle (\x+0.6,0.1);}
  \draw[arr,accent,line width=1.2pt] (-3.7,0) -- (-2.9,0);
  \draw[arr,accent,line width=1.2pt] (3.7,0) -- (2.9,0);
  \node[lbl] at (-2.8,-0.5) {левый горизонтальный};
  \node[lbl] at (2.8,-0.5) {правый горизонтальный};
  \draw[fill=gray!30] (0,-0.6) circle (0.2);
  \node[lbl] at (0,-0.95) {питание};
  \node[font=\sffamily\small\bfseries,text=accent] at (0,0.6) {триммеры внутрь + питание};
\end{tikzpicture}
\caption{Вход в загрузчик EdgeTX}
\end{figure}

Загрузчик тоже иногда обновляют: после обновления прошивки откройте \mpath{SYS\sep SD card},
найдите файл прошивки в \code{FIRMWARE}, долгим \btn{ENTER} вызовите меню и выберите
\menu{Flash bootloader}.

\section{Обновление внутреннего модуля ExpressLRS}
\label{sec:elrs-update}

Внутренний модуль ELRS --- это отдельный микроконтроллер (ESP32) со своей прошивкой.
Её версия должна быть согласована с версией приёмников: \textbf{мажорные версии передатчика и
приёмника должны совпадать} (например, 3.x с 3.x), иначе привязка не получится. Перед
переходом на новую мажорную версию прочитайте заметки к выпуску ExpressLRS.

\subsection{Через Wi-Fi (самый простой способ)}

\begin{enumerate}
  \item На компьютере установите \textbf{ExpressLRS Configurator}, выберите версию,
        категорию \menu{RadioMaster 2.4 GHz} и цель (target) вашего пульта: например,
        \menu{RadioMaster TX12 2400 TX} или \menu{RadioMaster Boxer 2400 TX}
        (точные названия зависят от версии Configurator).
  \item Задайте \textbf{фразу привязки} (binding phrase) --- любую строку, одинаковую для
        передатчика и всех ваших приёмников. Метод прошивки --- \menu{WiFi}. Нажмите
        \menu{Build}: Configurator соберёт файл \code{firmware.bin}.
  \item На пульте откройте \mpath{SYS\sep Tools\sep ExpressLRS} и выберите
        \menu{WiFi Connectivity\sep Enable WiFi}.
  \item Подключите компьютер к сети \code{ExpressLRS TX} (пароль \code{expresslrs}) и
        откройте в браузере \code{http://10.0.0.1}.
  \item На странице модуля загрузите \code{firmware.bin} и дождитесь перезагрузки модуля.
\end{enumerate}

\subsection{Через USB (EdgeTX Passthrough)}

В Configurator выберите метод \menu{EdgeTX Passthrough}. Пульт включите и подключите по USB,
выбрав режим \menu{USB Serial (VCP)}. Configurator передаст прошивку в модуль «сквозь» пульт.

\section{Обновление мультипротокольного модуля}

Для версий с 4-в-1 модулем:
\begin{enumerate}
  \item Скачайте прошивку Multi (\code{.bin}) для STM32 с сайта проекта
        DIY-Multiprotocol-TX-Module (обычно вариант с поддержкой всех протоколов и
        телеметрии для EdgeTX/OpenTX).
  \item Положите файл в папку \code{FIRMWARE} на SD-карте.
  \item \mpath{SYS\sep SD card}, найдите файл, долгое \btn{ENTER} → \menu{Flash internal Multi}.
\end{enumerate}

\begin{summary}
\begin{itemize}
  \item Перед любым обновлением --- полная копия SD-карты.
  \item Прошивка пульта: EdgeTX Buddy (DFU) или загрузчик (триммеры внутрь + питание).
  \item После прошивки обновите содержимое SD-карты под ту же версию.
  \item Модуль ELRS прошивается отдельно, версия должна совпадать с приёмниками.
\end{itemize}
\end{summary}
